\renewcommand{\thesection}{\arabic{section}}\renewcommand{\theHsection}{chapter.\arabic{section}}\setcounter{section}{6}\section{无穷级数}
\noindent\begin{minipage}{\textwidth}\ProblemHeading{习题}{7.1}
证明下列数项级数的基本性质：
\SubQ{1}{Cauchy 收敛准则：$\sum_{n=1}^{\infty}a_n$ 收敛当且仅当，对任意 $\varepsilon>0$，存在 $N$，使 $n>N$、$p\ge1$ 时 $|a_{n+1}+\cdots+a_{n+p}|<\varepsilon$。}
\SubQ{2}{若 $\sum a_n$ 收敛，则 $a_n\to0$。}
\SubQ{3}{若 $\sum a_n=A$、$\sum b_n=B$ 收敛，$\alpha$ 为常数，则 $\sum\alpha a_n=\alpha A$，$\sum(a_n\pm b_n)=A\pm B$。}
\end{minipage}\par\vfill
\noindent\begin{minipage}{\textwidth}\ProblemHeading{习题}{7.2}
证明下列数项级数的基本性质：
\SubQ{1}{改变级数的任意有限项不影响其敛散性。}
\end{minipage}\par\vfill
\noindent\begin{minipage}{\textwidth}\ProblemHeading{例题}{7.3}
证明 $\displaystyle\sum_{n=0}^{\infty}\frac1{n!}$ 收敛。
\end{minipage}\par\vfill
\noindent\begin{minipage}{\textwidth}\ProblemHeading{例题}{7.4}
证明调和级数 $\displaystyle\sum_{n=1}^{\infty}\frac1n$ 发散。
\end{minipage}\par\vfill
\newpage
\noindent\begin{minipage}{\textwidth}\ProblemHeading{例题}{7.5}
讨论 $p$ 级数 $\displaystyle\sum_{n=1}^{\infty}\frac1{n^p}$ 的敛散性，其中 $p\in\mathbb R$。
\end{minipage}\par\vfill
\noindent\begin{minipage}{\textwidth}\ProblemHeading{例题}{7.6}
证明 $\displaystyle\sum_{n=2}^{\infty}\frac1{\ln n}$ 发散。
\end{minipage}\par\vfill
\noindent\begin{minipage}{\textwidth}\ProblemHeading{例题}{7.7}
证明 $\displaystyle\sum_{n=1}^{\infty}\frac{n+3}{\sqrt{(n^2+1)(n^3+2)}}$ 收敛。
\end{minipage}\par\vfill
\noindent\begin{minipage}{\textwidth}\ProblemHeading{例题}{7.8}
证明 $\displaystyle\sum_{n=1}^{\infty}\sin\frac1n$ 发散。
\end{minipage}\par\vfill
\newpage
\noindent\begin{minipage}{\textwidth}\ProblemHeading{习题}{7.9}
设 $a_n\ge0$ 且 $\lim_{n\to\infty}\sqrt[n]{a_n}=q$。证明：$q<1$ 时 $\sum a_n$ 收敛，$q>1$ 时发散；并说明 $q=1$ 时不能据此判断敛散性。
\end{minipage}\par\vfill
\noindent\begin{minipage}{\textwidth}\ProblemHeading{例题}{7.10}
证明 $\displaystyle\sum_{n=1}^{\infty}\frac1{2^n}\left(1+\frac1n\right)^{n^2}$ 发散。
\end{minipage}\par\vfill
\noindent\begin{minipage}{\textwidth}\ProblemHeading{例题}{7.11}
讨论 $\displaystyle\sum_{n=1}^{\infty}n!\left(\frac xn\right)^n$（$x\ge0$）的敛散性，包括临界值 $x=e$。
\end{minipage}\par\vfill
\noindent\begin{minipage}{\textwidth}\ProblemHeading{例题}{7.12}
证明 $\displaystyle\sum_{n=2}^{\infty}\frac1{n\ln^\alpha n}$ 在 $\alpha>1$ 时收敛、$\alpha\le1$ 时发散。
\end{minipage}\par\vfill
\newpage
\noindent\begin{minipage}{\textwidth}\ProblemHeading{习题}{7.13}
证明 Raabe 判别法：若 $a_n>0$，$\lim_{n\to\infty}n\left(\frac{a_n}{a_{n+1}}-1\right)=\alpha$，则 $\alpha>1$ 时 $\sum a_n$ 收敛，$0<\alpha<1$ 时发散。
\end{minipage}\par\vfill
\noindent\begin{minipage}{\textwidth}\ProblemHeading{例题}{7.14（改写）}
设 $\alpha,\beta,\gamma,x>0$，讨论超几何级数\[1+\sum_{n=1}^{\infty}\frac{\alpha(\alpha+1)\cdots(\alpha+n-1)\beta(\beta+1)\cdots(\beta+n-1)}{n!\,\gamma(\gamma+1)\cdots(\gamma+n-1)}x^n\]的敛散性。
\end{minipage}\par\vfill
\noindent\begin{minipage}{\textwidth}\ProblemHeading{例题}{7.15（改写）}
设 $\sum_{n=1}^{\infty}a_n$ 的部分和有界，$\alpha>0$。证明 $\sum_{n=1}^{\infty}a_n/n^\alpha$ 收敛。
\end{minipage}\par\vfill
\noindent\begin{minipage}{\textwidth}\ProblemHeading{例题}{7.16}
讨论级数 $\displaystyle\sum_{n=1}^{\infty}\frac{\cos nx}{n}$ 的敛散性，其中 $x\in\mathbb R$。
\end{minipage}\par\vfill
\newpage
\noindent\begin{minipage}{\textwidth}\ProblemHeading{例题}{7.17}
证明 $\displaystyle\sum_{n=1}^{\infty}\frac{\cos n}{n}$ 条件收敛。
\end{minipage}\par\vfill
\noindent\begin{minipage}{\textwidth}\ProblemHeading{例题}{7.18（改写）}
设 $b_n$ 单调递减且趋于零。利用 Dirichlet 判别法证明交错级数 $\sum_{n=1}^{\infty}(-1)^{n+1}b_n$ 收敛。
\end{minipage}\par\vfill
\noindent\begin{minipage}{\textwidth}\ProblemHeading{例题}{7.19（改写）}
证明 $\displaystyle\sum_{n=1}^{\infty}\frac{(-1)^{n-1}}{\sqrt n}$ 收敛，但它与自身的 Cauchy 乘积发散。这里两个级数的 Cauchy 乘积定义为 $\sum_{n=1}^{\infty}c_n$，其中 $c_n=\sum_{i+j=n+1}a_i b_j$。
\end{minipage}\par\vfill
\noindent\begin{minipage}{\textwidth}\ProblemHeading{习题}{7.20}
设 $\sum_{n=1}^{\infty}a_n=A$、$\sum_{n=1}^{\infty}b_n=B$ 都绝对收敛。将全部乘积 $a_i b_j$（$i,j\ge1$）按任意顺序各取一次排列为一个级数。证明所得级数绝对收敛，且和为 $AB$。
\end{minipage}\par\vfill
\newpage
\noindent\begin{minipage}{\textwidth}\ProblemHeading{习题}{7.21}
证明下列等式：
\SubQ{1}{$\displaystyle\sum_{n=1}^{\infty}\frac1{(2n-1)(2n+1)}=\frac12$。}
\SubQ{2}{$\displaystyle\sum_{n=1}^{\infty}(\sqrt{n+2}-2\sqrt{n+1}+\sqrt n)=1-\sqrt2$。}
\SubQ{3}{$\displaystyle\sum_{n=1}^{\infty}\ln\frac{n(2n+1)}{(n+1)(2n-1)}=\ln2$。}
\end{minipage}\par\vfill
\noindent\begin{minipage}{\textwidth}\ProblemHeading{习题}{7.22}
证明下列等式：
\SubQ{1}{$\displaystyle\sum_{n=1}^{\infty}\frac{2n+1}{n^2(n+1)^2}=1$。}
\end{minipage}\par\vfill
\noindent\begin{minipage}{\textwidth}\ProblemHeading{习题}{7.23}
研究下列级数的敛散性，其中未另作说明的参数为实数：
\SubQ{1}{$\displaystyle\sum_{n=1}^{\infty}\sqrt[n]{0.001}$。}
\SubQ{2}{$\displaystyle\sum_{n=2}^{\infty}\frac1{n\sqrt{n-1}}$。}
\SubQ{3}{$\displaystyle\sum_{n=1}^{\infty}\frac1{\sqrt{(2n-1)(2n+1)}}$。}
\end{minipage}\par\vfill
\noindent\begin{minipage}{\textwidth}\ProblemHeading{习题}{7.24}
研究下列级数的敛散性，其中未另作说明的参数为实数：
\SubQ{1}{$\displaystyle\sum_{n=1}^{\infty}\sin n$。}
\SubQ{2}{$\displaystyle\sum_{n=1}^{\infty}2^n\sin\frac\pi{3^n}$。}
\SubQ{3}{$\displaystyle\sum_{n=1}^{\infty}\frac1{n\sqrt[n]n}$。}
\end{minipage}\par\vfill
\newpage
\noindent\begin{minipage}{\textwidth}\ProblemHeading{习题}{7.25}
研究下列级数的敛散性，其中未另作说明的参数为实数：
\SubQ{1}{$\displaystyle\sum_{n=1}^{\infty}\frac1{(2+1/n)^n}$。}
\SubQ{2}{$\displaystyle\sum_{n=1}^{\infty}\frac n{(n+1/n)^n}$。}
\SubQ{3}{$\displaystyle\sum_{n=1}^{\infty}\arctan\frac\pi{4n}$。}
\end{minipage}\par\vfill
\noindent\begin{minipage}{\textwidth}\ProblemHeading{习题}{7.26}
研究下列级数的敛散性，其中未另作说明的参数为实数：
\SubQ{1}{$\displaystyle\sum_{n=1}^{\infty}\frac{1000^n}{n!}$。}
\SubQ{2}{$\displaystyle\sum_{n=1}^{\infty}\frac{(n!)^2}{(2n)!}$。}
\SubQ{3}{$\displaystyle\sum_{n=1}^{\infty}\frac{3+(-1)^n}{2^n}$。}
\end{minipage}\par\vfill
\noindent\begin{minipage}{\textwidth}\ProblemHeading{习题}{7.27}
研究下列级数的敛散性，其中未另作说明的参数为实数：
\SubQ{1}{$\displaystyle\sum_{n=1}^{\infty}\frac{\ln n}{\sqrt[4]{n^5}}$。}
\SubQ{2}{$\displaystyle\sum_{n=3}^{\infty}\frac1{n\ln n(\ln\ln n)^k}$。}
\SubQ{3}{$\displaystyle\sum_{n=1}^{\infty}\left(\cos\frac1n\right)^{n^3}$。}
\end{minipage}\par\vfill
\noindent\begin{minipage}{\textwidth}\ProblemHeading{习题}{7.28}
研究下列级数的敛散性，其中未另作说明的参数为实数：
\SubQ{1}{$\displaystyle\sum_{n=2}^{\infty}\left(\frac{an}{n+1}\right)^n$，$a>0$。}
\end{minipage}\par\vfill
\newpage
\noindent\begin{minipage}{\textwidth}\ProblemHeading{习题}{7.29}
证明并举例：
\SubQ{1}{若 $\sum_{n=1}^{\infty}a_n$ 收敛，则 $\sum_{n=1}^{\infty}(a_n+a_{n+1})$ 收敛。}
\SubQ{2}{举例说明上述逆命题不成立。}
\SubQ{3}{若再要求 $a_n>0$，证明上述逆命题成立。}
\end{minipage}\par\vfill
\noindent\begin{minipage}{\textwidth}\ProblemHeading{习题}{7.30}
证明或回答下列论断：
\SubQ{1}{若 $\lim_{n\to\infty}na_n=a\ne0$，则 $\sum_{n=1}^{\infty}a_n$ 发散。}
\SubQ{2}{若 $\sum_{n=1}^{\infty}a_n$ 收敛，是否必有 $\lim_{n\to\infty}na_n=0$？}
\SubQ{3}{若 $\lim_{n\to\infty}na_n=a$ 为有限实数，且 $\sum_{n=1}^{\infty}n(a_n-a_{n+1})$ 收敛，则 $\sum_{n=1}^{\infty}a_n$ 收敛。}
\end{minipage}\par\vfill
\noindent\begin{minipage}{\textwidth}\ProblemHeading{习题}{7.31}
设正项级数 $\sum a_n$ 收敛。证明 $\sum a_n^2$ 也收敛，并判断逆命题是否成立。
\end{minipage}\par\vfill
\noindent\begin{minipage}{\textwidth}\ProblemHeading{习题}{7.32}
设 $a_n,b_n\ge0$，$a_{n+1}<a_n+b_n$，且 $\sum_{n=1}^{\infty}b_n$ 收敛。证明 $\lim_{n\to\infty}a_n$ 存在。
\end{minipage}\par\vfill
\newpage
\noindent\begin{minipage}{\textwidth}\ProblemHeading{习题}{7.33}
设 $\sum a_n^2$ 和 $\sum b_n^2$ 收敛。证明下列级数均收敛：
\SubQ{1}{$\sum_{n=1}^{\infty}|a_nb_n|$。}
\SubQ{2}{$\sum_{n=1}^{\infty}(a_n+b_n)^2$。}
\SubQ{3}{$\sum_{n=1}^{\infty}|a_n|/n$。}
\end{minipage}\par\vfill
\noindent\begin{minipage}{\textwidth}\ProblemHeading{习题}{7.34}
求下列极限，其中 $p>1$：
\SubQ{1}{$\displaystyle\lim_{n\to\infty}\left(\frac1{(n+1)^p}+\frac1{(n+2)^p}+\cdots+\frac1{(2n)^p}\right)$。}
\SubQ{2}{$\displaystyle\lim_{n\to\infty}\left(\frac1{p^{n+1}}+\frac1{p^{n+2}}+\cdots+\frac1{p^{2n}}\right)$。}
\end{minipage}\par\vfill
\noindent\begin{minipage}{\textwidth}\ProblemHeading{习题}{7.35}
设正项数列 $a_n$ 单调递减，且 $\sum_{n=1}^{\infty}(-1)^n a_n$ 发散。判断 $\sum_{n=1}^{\infty}(1/(a_n+1))^n$ 是否收敛，并说明理由。
\end{minipage}\par\vfill
\noindent\begin{minipage}{\textwidth}\ProblemHeading{习题}{7.36}
设 $a_n>a_{n+1}>0$ 且 $a_n\to0$。证明 $\displaystyle\sum_{n=1}^{\infty}(-1)^{n-1}\frac{a_1+\cdots+a_n}{n}$ 收敛。
\end{minipage}\par\vfill
\newpage
\noindent\begin{minipage}{\textwidth}\ProblemHeading{习题}{7.37}
设 $\sum a_n$ 和 $\sum b_n$ 绝对收敛。证明 $\sum(a_n+b_n)$ 绝对收敛。
\end{minipage}\par\vfill
\noindent\begin{minipage}{\textwidth}\ProblemHeading{习题}{7.38}
研究下列级数的条件收敛性与绝对收敛性，其中 $p\in\mathbb R$：
\SubQ{1}{$\displaystyle\sum_{n=1}^{\infty}(-1)^n\left(\frac{2n+100}{3n+1}\right)^n$。}
\SubQ{2}{$\displaystyle\sum_{n=1}^{\infty}\frac{(-1)^{n(n-1)/2}}{2^n}$。}
\SubQ{3}{$\displaystyle\sum_{n=1}^{\infty}(-1)^n\frac{\sqrt n}{n+100}$。}
\end{minipage}\par\vfill
\noindent\begin{minipage}{\textwidth}\ProblemHeading{习题}{7.39}
研究下列级数的条件收敛性与绝对收敛性，其中 $p\in\mathbb R$：
\SubQ{1}{$\displaystyle\sum_{n=1}^{\infty}(-1)^{n-1}\sin\frac1n$。}
\SubQ{2}{$\displaystyle\sum_{n=1}^{\infty}(-1)^{n-1}\frac{\ln n}{n}$。}
\SubQ{3}{$\displaystyle\sum_{n=1}^{\infty}\frac{(-1)^{n-1}}{n^p}$。}
\end{minipage}\par\vfill
\noindent\begin{minipage}{\textwidth}\ProblemHeading{习题}{7.40}
研究下列级数的条件收敛性与绝对收敛性，其中 $p\in\mathbb R$：
\SubQ{1}{$\displaystyle\sum_{n=1}^{\infty}(-1)^n(e^{1/n}-1)$。}
\SubQ{2}{$\displaystyle\sum_{n=1}^{\infty}(-1)^n\left(\frac1n-\ln\left(1+\frac1n\right)\right)$。}
\SubQ{3}{$\displaystyle\sum_{n=1}^{\infty}(-1)^n\left(1-\cos\frac pn\right)$。}
\end{minipage}\par\vfill
\newpage
\noindent\begin{minipage}{\textwidth}\ProblemHeading{习题}{7.41}
研究下列级数的条件收敛性与绝对收敛性，其中 $p\in\mathbb R$：
\SubQ{1}{$\displaystyle\sum_{n=1}^{\infty}(-1)^n\left(1-\cos\frac1n\right)^p$。}
\end{minipage}\par\vfill
\noindent\begin{minipage}{\textwidth}\ProblemHeading{习题}{7.42}
设 $\sum a_n$ 条件收敛，令 $a_n^+=(|a_n|+a_n)/2$、$a_n^-=(|a_n|-a_n)/2$，$S_n^\pm=\sum_{k=1}^n a_k^\pm$。证明 $\lim_{n\to\infty}S_n^+/S_n^-=1$。
\end{minipage}\par\vfill
\noindent\begin{minipage}{\textwidth}\ProblemHeading{习题}{7.43}
设正项级数 $\sum b_n$ 收敛，且从某项起 $a_n\ne0$ 并有 $|a_{n+1}/a_n|<b_{n+1}/b_n$。证明 $\sum a_n$ 绝对收敛。
\end{minipage}\par\vfill
\noindent\begin{minipage}{\textwidth}\ProblemHeading{习题}{7.44}
研究下列级数的敛散性，其中 $x\in\mathbb R$：
\SubQ{1}{$\displaystyle\sum_{n=1}^{\infty}\frac{\sin nx}{n}$。}
\SubQ{2}{$\displaystyle\sum_{n=2}^{\infty}\frac{\cos(n\pi/4)}{\ln n}$。}
\SubQ{3}{$\displaystyle\sum_{n=1}^{\infty}\frac{\sin n}{\sqrt n}\left(1+\frac1n\right)^n$。}
\end{minipage}\par\vfill
\newpage
\noindent\begin{minipage}{\textwidth}\ProblemHeading{习题}{7.45}
研究下列级数的敛散性，其中 $x\in\mathbb R$：
\SubQ{1}{$\displaystyle\sum_{n=1}^{\infty}(-1)^n\frac{n-1}{n+1}\frac1{\sqrt[100]n}$。}
\end{minipage}\par\vfill
\noindent\begin{minipage}{\textwidth}\ProblemHeading{例题}{7.46}
讨论 $\displaystyle\sum_{n=0}^{\infty}x^n$ 的收敛性。
\end{minipage}\par\vfill
\noindent\begin{minipage}{\textwidth}\ProblemHeading{例题}{7.47（改写）}
在 $[0,1]$ 上考虑级数 $x+x(x-1)+x^2(x-1)+\cdots$。求其部分和与和函数，说明和函数的连续性。
\end{minipage}\par\vfill
\noindent\begin{minipage}{\textwidth}\ProblemHeading{例题}{7.48（改写）}
设 $r_1,r_2,\ldots$ 是 $[0,1]$ 中所有有理数的一个枚举，并定义\[S_n(x)=\begin{cases}1,&x\in\{r_1,\ldots,r_n\},\\0,&\text{其他情形}.\end{cases}\]求逐点极限 $S(x)$，比较 $S_n$ 与 $S$ 在 $[0,1]$ 上的 Riemann 可积性。
\end{minipage}\par\vfill
\newpage
\noindent\begin{minipage}{\textwidth}\ProblemHeading{例题}{7.49（改写）}
设 $S_n(x)=2n^2xe^{-n^2x^2}$（$x\in[0,1]$）。求逐点极限 $S(x)$，分别计算 $\int_0^1 S(x)\,dx$ 和 $\lim_{n\to\infty}\int_0^1 S_n(x)\,dx$，比较二者。
\end{minipage}\par\vfill
\noindent\begin{minipage}{\textwidth}\ProblemHeading{例题}{7.50（改写）}
设 $S_n(x)=\sin(nx)/\sqrt n$（$x\in\mathbb R$）。求逐点极限 $S(x)$；计算 $S_n^{\prime}(x)$，说明逐项取极限与求导是否可以交换。
\end{minipage}\par\vfill
\noindent\begin{minipage}{\textwidth}\ProblemHeading{习题}{7.51}
证明：函数列 $f_n$ 在区间 $I$ 上一致收敛于 $f$ 的充分必要条件是 $\lim_{n\to\infty}\sup_{x\in I}|f_n(x)-f(x)|=0$。
\end{minipage}\par\vfill
\noindent\begin{minipage}{\textwidth}\ProblemHeading{例题}{7.52}
讨论函数列 $f_n(x)=1/(x+n)$ 在 $[0,1]$ 上的一致收敛性。
\end{minipage}\par\vfill
\newpage
\noindent\begin{minipage}{\textwidth}\ProblemHeading{例题}{7.53（改写）}
研究函数列 $S_n(x)=x^n$ 在 $[0,1]$ 上的逐点极限及一致收敛性，并据此判断级数 $x+x(x-1)+x^2(x-1)+\cdots$ 是否一致收敛。
\end{minipage}\par\vfill
\noindent\begin{minipage}{\textwidth}\ProblemHeading{例题}{7.54}
讨论 $\displaystyle\sum_{n=1}^{\infty}ne^{-nx}$ 在 $(0,+\infty)$ 上的一致收敛性。
\end{minipage}\par\vfill
\noindent\begin{minipage}{\textwidth}\ProblemHeading{例题}{7.55}
设 $\alpha>1$。证明 $\displaystyle\sum_{n=1}^{\infty}\frac{\cos nx}{n^\alpha}$ 在 $\mathbb R$ 上一致收敛。
\end{minipage}\par\vfill
\noindent\begin{minipage}{\textwidth}\ProblemHeading{习题}{7.56}
证明函数项级数的一致收敛判别法，并说明证明中各步的一致性。设所有函数均在区间 $E$ 上定义：
\SubQ{1}{Dirichlet 判别法：对每个 $x\in E$，$v_n(x)$ 关于 $n$ 单调递减，且 $v_n$ 在 $E$ 上一致趋于零；若 $\sum u_n(x)$ 的部分和在 $E$ 上一致有界，则 $\sum u_n(x)v_n(x)$ 在 $E$ 上一致收敛。}
\SubQ{2}{Abel 判别法：对每个 $x\in E$，$v_n(x)$ 关于 $n$ 单调，且函数列 $v_n$ 在 $E$ 上一致有界；若 $\sum u_n(x)$ 在 $E$ 上一致收敛，则 $\sum u_n(x)v_n(x)$ 在 $E$ 上一致收敛。}
\end{minipage}\par\vfill
\newpage
\noindent\begin{minipage}{\textwidth}\ProblemHeading{例题}{7.57}
设 $a_n$ 单调递减趋于零。证明对任意 $0<\delta<\pi$，$\sum_{n=1}^{\infty}a_n\cos nx$ 在 $[\delta,2\pi-\delta]$ 上一致收敛。
\end{minipage}\par\vfill
\noindent\begin{minipage}{\textwidth}\ProblemHeading{例题}{7.58}
设 $\sum_{n=1}^{\infty}a_n$ 收敛。证明 $\sum_{n=1}^{\infty}a_n/n^x$ 在 $[0,+\infty)$ 上一致收敛。
\end{minipage}\par\vfill
\noindent\begin{minipage}{\textwidth}\ProblemHeading{例题}{7.59}
设 $f(x)=\sum_{n=1}^{\infty}\cos(nx)/n^2$。求 $\int_0^\pi f(x)\,dx$。
\end{minipage}\par\vfill
\noindent\begin{minipage}{\textwidth}\ProblemHeading{习题}{7.60}
证明：两个在同一区间 $I$ 上一致收敛的级数的逐项和也在 $I$ 上一致收敛。
\end{minipage}\par\vfill
\newpage
\noindent\begin{minipage}{\textwidth}\ProblemHeading{习题}{7.61}
确定下列函数项级数的收敛域（$x$ 取实数）：
\SubQ{1}{$\displaystyle\sum_{n=1}^{\infty}ne^{-nx}$。}
\SubQ{2}{$\displaystyle\sum_{n=2}^{\infty}\frac{x^{n^2}}n$。}
\SubQ{3}{$\displaystyle\sum_{n=1}^{\infty}\frac{(-1)^n}{2n-1}\left(\frac{1-x}{1+x}\right)^n$。}
\end{minipage}\par\vfill
\noindent\begin{minipage}{\textwidth}\ProblemHeading{习题}{7.62}
确定下列函数项级数的收敛域（$x$ 取实数）：
\SubQ{1}{$\displaystyle\sum_{n=1}^{\infty}\frac1{x^n}\sin\frac\pi{2^n}$。}
\SubQ{2}{$\displaystyle\sum_{n=1}^{\infty}\frac{(x-3)^n}{n-3^n}$。}
\SubQ{3}{$\displaystyle\sum_{n=1}^{\infty}n!\left(\frac xn\right)^n$。}
\end{minipage}\par\vfill
\noindent\begin{minipage}{\textwidth}\ProblemHeading{习题}{7.63}
确定下列函数项级数的收敛域（$x$ 取实数）：
\SubQ{1}{$\displaystyle\sum_{n=1}^{\infty}\frac{\cos nx}{e^{nx}}$。}
\SubQ{2}{$\displaystyle\sum_{n=1}^{\infty}\frac{x^n}{1-x^n}$。}
\end{minipage}\par\vfill
\noindent\begin{minipage}{\textwidth}\ProblemHeading{习题}{7.64}
在 $[0,1]$ 上定义\[u_n(x)=\begin{cases}1/n,&x=1/n,\\0,&x\ne1/n.\end{cases}\]证明 $\sum_{n=1}^{\infty}u_n(x)$ 一致收敛，但不存在收敛正项级数 $\sum M_n$ 使所有 $x\in[0,1]$ 都满足 $|u_n(x)|\le M_n$。
\end{minipage}\par\vfill
\newpage
\noindent\begin{minipage}{\textwidth}\ProblemHeading{习题}{7.65}
研究下列级数在给定区间上的一致收敛性：
\SubQ{1}{$\displaystyle\sum_{n=1}^{\infty}\frac{\sin nx}{n^2}$，$x\in\mathbb R$。}
\SubQ{2}{$\displaystyle\sum_{n=1}^{\infty}\frac1{2^n[1+(nx)^2]}$，$x\in\mathbb R$。}
\SubQ{3}{$\displaystyle\sum_{n=1}^{\infty}(-1)^{n-1}x^n$，$-1<x<1$。}
\end{minipage}\par\vfill
\noindent\begin{minipage}{\textwidth}\ProblemHeading{习题}{7.66}
研究下列级数在给定区间上的一致收敛性：
\SubQ{1}{$\displaystyle\sum_{n=1}^{\infty}x^2e^{-nx}$，$0\le x<+\infty$。}
\SubQ{2}{$\displaystyle\sum_{n=1}^{\infty}\frac{(-1)^n}{x+n}$，$0\le x<+\infty$。}
\SubQ{3}{$\displaystyle\sum_{n=1}^{\infty}\frac1{n^x}$，$1<x<+\infty$。}
\end{minipage}\par\vfill
\noindent\begin{minipage}{\textwidth}\ProblemHeading{习题}{7.67}
研究下列级数在给定区间上的一致收敛性：
\SubQ{1}{$\displaystyle\sum_{n=1}^{\infty}\frac{\cos nx}{n}$，$\delta\le x\le2\pi-\delta$，$0<\delta<\pi$。}
\SubQ{2}{$\displaystyle\sum_{n=1}^{\infty}\frac{x^2}{(ne^n)^x}$，$0\le x<+\infty$。}
\end{minipage}\par\vfill
\noindent\begin{minipage}{\textwidth}\ProblemHeading{习题}{7.68}
证明：若 $\sum a_n$ 收敛，则 $\sum_{n=1}^{\infty}a_n e^{-nx}$ 在 $[0,+\infty)$ 上一致收敛。
\end{minipage}\par\vfill
\newpage
\noindent\begin{minipage}{\textwidth}\ProblemHeading{习题}{7.69}
证明 $\zeta(x)=\sum_{n=1}^{\infty}1/n^x$ 在 $(1,+\infty)$ 内连续，且具有连续的各阶导数。
\end{minipage}\par\vfill
\noindent\begin{minipage}{\textwidth}\ProblemHeading{习题}{7.70}
证明 $f(x)=\sum_{n=1}^{\infty}\sin(nx)/n^4$ 在实轴上具有连续的二阶导数。
\end{minipage}\par\vfill
\noindent\begin{minipage}{\textwidth}\ProblemHeading{习题}{7.71}
设 $\displaystyle f(x)=\sum_{n=1}^{\infty}\frac{x^n\cos(n\pi/x)}{(1+2x)^n}$。求 $\lim_{x\to+\infty}f(x)$ 及 $\lim_{x\to1}f(x)$。
\end{minipage}\par\vfill
\noindent\begin{minipage}{\textwidth}\ProblemHeading{习题}{7.72}
设 $f(x)=\sum_{n=1}^{\infty}ne^{-nx}$。求 $\int_{\ln2}^{\ln3}f(x)\,dx$。
\end{minipage}\par\vfill
\newpage
\noindent\begin{minipage}{\textwidth}\ProblemHeading{习题}{7.73}
在 $[0,1)$ 上递归定义连续可微函数列 $f_n$：$f_1=1$，$f_{n+1}^{\prime}(x)=f_n(x)f_{n+1}(x)$（$0<x<1$），$f_{n+1}(0)=1$。证明对每个 $x\in[0,1)$，$\lim_{n\to\infty}f_n(x)$ 存在，并求极限函数。
\end{minipage}\par\vfill
\noindent\begin{minipage}{\textwidth}\ProblemHeading{习题}{7.74}
证明 Dini 定理的下列形式：
\SubQ{1}{设 $u_n$ 是 $[a,b]$ 上的非负连续函数，逐点趋于零；若对每个固定 $x$，$u_n(x)$ 关于 $n$ 单调递减，则 $u_n$ 在 $[a,b]$ 上一致趋于零。}
\SubQ{2}{若 $\sum u_n(x)$ 在 $[a,b]$ 上逐点收敛于 $S(x)$，每个 $u_n$ 连续且非负，则 $S$ 连续当且仅当该级数一致收敛。}
\end{minipage}\par\vfill
\noindent\begin{minipage}{\textwidth}\ProblemHeading{例题}{7.75}
求下列幂级数的收敛区域：
\SubQ{1}{$1+\sum_{n=1}^{\infty}n^n x^n$。}
\SubQ{2}{$\sum_{n=0}^{\infty}x^n/n!$。}
\end{minipage}\par\vfill
\noindent\begin{minipage}{\textwidth}\ProblemHeading{例题}{7.76}
求级数 $\sum_{n=1}^{\infty}n^\alpha x^n$ 的收敛半径与收敛区域，其中 $\alpha\in\mathbb R$。
\end{minipage}\par\vfill
\newpage
\noindent\begin{minipage}{\textwidth}\ProblemHeading{例题}{7.77}
求级数 $\displaystyle\sum_{n=1}^{\infty}\frac{(2n)!}{(n!)^2}x^{2n}$ 的收敛半径。
\end{minipage}\par\vfill
\noindent\begin{minipage}{\textwidth}\ProblemHeading{例题}{7.78}
求 $\sum_{n=0}^{\infty}a_nx^n$ 的收敛半径，其中 $a_{2k}=2^k$，$a_{2k+1}=k$（$k=0,1,2,\ldots$）。
\end{minipage}\par\vfill
\noindent\begin{minipage}{\textwidth}\ProblemHeading{例题}{7.79（改写）}
已知 $S(x)=\sum_{n=0}^{\infty}x^n/n!$ 在整个实轴上收敛。利用逐项求导及微分方程，求其和函数。
\end{minipage}\par\vfill
\noindent\begin{minipage}{\textwidth}\ProblemHeading{例题}{7.80}
求级数 $\sum_{n=0}^{\infty}nx^n$ 的和函数。
\end{minipage}\par\vfill
\newpage
\noindent\begin{minipage}{\textwidth}\ProblemHeading{例题}{7.81}
将 $\displaystyle\frac1{(1-x)(2-x)}$ 展开为 Maclaurin 级数。
\end{minipage}\par\vfill
\noindent\begin{minipage}{\textwidth}\ProblemHeading{例题}{7.82}
在 $x=3$ 处将 $\ln x$ 展开为 Taylor 级数。
\end{minipage}\par\vfill
\noindent\begin{minipage}{\textwidth}\ProblemHeading{习题}{7.83}
求下列幂级数的收敛半径：
\SubQ{1}{$\displaystyle\sum_{n=1}^{\infty}(-1)^{n+1}\frac{x^n}{n^2}$。}
\SubQ{2}{$\displaystyle\sum_{n=1}^{\infty}\frac{(n!)^2}{(2n)!}x^n$。}
\SubQ{3}{$\displaystyle\sum_{n=1}^{\infty}2^n x^{2n}$。}
\end{minipage}\par\vfill
\noindent\begin{minipage}{\textwidth}\ProblemHeading{习题}{7.84}
求下列幂级数的收敛半径：
\SubQ{1}{$\displaystyle\sum_{n=1}^{\infty}\frac{x^n}{a^n+b^n}$，$a,b>0$。}
\SubQ{2}{$\displaystyle\sum_{n=1}^{\infty}\frac{(x-2)^{2n-1}}{(2n-1)!}$。}
\SubQ{3}{$\displaystyle\sum_{n=1}^{\infty}\frac{3^n+(-2)^n}{n}(x+1)^n$。}
\end{minipage}\par\vfill
\newpage
\noindent\begin{minipage}{\textwidth}\ProblemHeading{习题}{7.85}
求下列幂级数的收敛半径：
\SubQ{1}{$\displaystyle\sum_{n=1}^{\infty}\left(1+\frac12+\cdots+\frac1n\right)x^n$。}
\SubQ{2}{$\displaystyle\sum_{n=1}^{\infty}\frac{x^{n^2}}{2^n}$。}
\end{minipage}\par\vfill
\noindent\begin{minipage}{\textwidth}\ProblemHeading{习题}{7.86}
设 $0<R<+\infty$，$f(x)=\sum_{n=0}^{\infty}a_nx^n$ 在 $|x|<R$ 时收敛，且 $\sum_{n=0}^{\infty}\frac{a_n}{n+1}R^{n+1}$ 收敛。
\SubQ{1}{证明 $\displaystyle\int_0^R f(x)\,dx=\sum_{n=0}^{\infty}\frac{a_n}{n+1}R^{n+1}$。积分在必要时按反常积分理解。}
\SubQ{2}{应用该结果证明 $\displaystyle\int_0^1\frac{dx}{1+x}=\ln2=\sum_{n=1}^{\infty}\frac{(-1)^{n-1}}n$。}
\end{minipage}\par\vfill
\noindent\begin{minipage}{\textwidth}\ProblemHeading{习题}{7.87}
求下列幂级数的收敛区域与和函数：
\SubQ{1}{$\displaystyle\sum_{n=0}^{\infty}(-1)^n\frac{x^{2n+1}}{2n+1}$。}
\SubQ{2}{$\displaystyle\sum_{n=0}^{\infty}(n+1)x^n$。}
\SubQ{3}{$\displaystyle\sum_{n=1}^{\infty}n(n+1)x^{n-1}$。}
\end{minipage}\par\vfill
\noindent\begin{minipage}{\textwidth}\ProblemHeading{习题}{7.88}
求下列幂级数的收敛区域与和函数：
\SubQ{1}{$\displaystyle\sum_{n=1}^{\infty}\frac{x^n}{n(n+1)}$。}
\SubQ{2}{$\displaystyle\sum_{n=1}^{\infty}\frac{x^{2n-1}}{(2n-1)!!}$，其中 $(2n-1)!!=1\cdot3\cdots(2n-1)$。}
\end{minipage}\par\vfill
\newpage
\noindent\begin{minipage}{\textwidth}\ProblemHeading{习题}{7.89}
求下列级数的和：
\SubQ{1}{$\displaystyle\sum_{n=2}^{\infty}\frac1{(n^2-1)2^n}$。}
\SubQ{2}{$\displaystyle\sum_{n=0}^{\infty}\frac{(-1)^n(n^2-n+1)}{2^n}$。}
\SubQ{3}{$\displaystyle\sum_{n=0}^{\infty}\frac{(-1)^n}{3n+1}$。}
\end{minipage}\par\vfill
\noindent\begin{minipage}{\textwidth}\ProblemHeading{习题}{7.90}
求下列级数的和：
\SubQ{1}{$\displaystyle\sum_{n=0}^{\infty}\frac{(n+1)^2}{n!}$。}
\end{minipage}\par\vfill
\noindent\begin{minipage}{\textwidth}\ProblemHeading{习题}{7.91}
求下列函数在指定点处的 Taylor 展开式，并给出收敛区域：
\SubQ{1}{$x^3-2x^2+5x-7$，$x_0=1$。}
\SubQ{2}{$e^{x/a}$，$x_0=a$（$a\ne0$）。}
\SubQ{3}{$\ln x$，$x_0=1$。}
\end{minipage}\par\vfill
\noindent\begin{minipage}{\textwidth}\ProblemHeading{习题}{7.92}
求下列函数在指定点处的 Taylor 展开式，并给出收敛区域：
\SubQ{1}{$1/(x^2+3x+2)$，$x_0=-4$。}
\SubQ{2}{$\ln(1+x-2x^2)$，$x_0=0$。}
\SubQ{3}{$\cos x$，$x_0=\pi/4$。}
\end{minipage}\par\vfill
\newpage
\noindent\begin{minipage}{\textwidth}\ProblemHeading{习题}{7.93}
求下列函数的 Maclaurin 展开式，并给出收敛区域：
\SubQ{1}{$\sin^2x$。}
\SubQ{2}{$\arcsin x$。}
\SubQ{3}{$\ln\sqrt{(1+x)/(1-x)}$。}
\end{minipage}\par\vfill
\noindent\begin{minipage}{\textwidth}\ProblemHeading{习题}{7.94}
求下列函数的 Maclaurin 展开式，并给出收敛区域：
\SubQ{1}{$(1+x)\ln(1+x)$。}
\SubQ{2}{$\int_0^x\cos(t^2)\,dt$。}
\SubQ{3}{$\int_0^x(\sin t)/t\,dt$（在 $t=0$ 连续延拓）。}
\end{minipage}\par\vfill
\noindent\begin{minipage}{\textwidth}\ProblemHeading{习题}{7.95}
求下列函数的 Maclaurin 展开式，并给出收敛区域：
\SubQ{1}{$\int_0^x e^{-t^2}\,dt$。}
\end{minipage}\par\vfill
\noindent\begin{minipage}{\textwidth}\ProblemHeading{习题}{7.96}
设 $\lambda\ne-1$。方程 $y+\lambda\sin y=x$ 在 $(x,y)=(0,0)$ 附近确定隐函数 $y(x)$。求它的幂级数展开式中 $x,x^3,x^5,x^7$ 各项的系数。
\end{minipage}\par\vfill
\newpage
\noindent\begin{minipage}{\textwidth}\ProblemHeading{例题}{7.97}
用级数方法计算 $\displaystyle\int_0^1\frac{\ln(1+x)}x\,dx$，在 $x=0$ 处按连续延拓理解。
\end{minipage}\par\vfill
\noindent\begin{minipage}{\textwidth}\ProblemHeading{例题}{7.98}
用级数表示椭圆积分 $\displaystyle\int_0^1\frac{du}{\sqrt{(1-u^2)(1-k^2u^2)}}$，其中 $k\in\mathbb R$、$k^2<1$。
\end{minipage}\par\vfill
\noindent\begin{minipage}{\textwidth}\ProblemHeading{例题}{7.99}
求 Airy 方程 $y^{\prime\prime}-xy=0$ 在 $x=0$ 附近的幂级数解。
\end{minipage}\par\vfill
\noindent\begin{minipage}{\textwidth}\ProblemHeading{例题}{7.100}
设 $\nu\ge0$，求解 Bessel 方程 $x^2y^{\prime\prime}+xy^{\prime}+(x^2-\nu^2)y=0$（$x>0$）。
\end{minipage}\par\vfill
\newpage
\noindent\begin{minipage}{\textwidth}\ProblemHeading{习题}{7.101}
用级数方法求下列积分：
\SubQ{1}{$\int_0^1e^{-x^2}\,dx$。}
\SubQ{2}{$\int_0^1(\sin x)/x\,dx$（在 $x=0$ 连续延拓）。}
\end{minipage}\par\vfill
\noindent\begin{minipage}{\textwidth}\ProblemHeading{习题}{7.102}
求方程 $y^{\prime\prime}-xy^{\prime}+y=0$ 在 $x=0$ 附近的幂级数解。
\end{minipage}\par\vfill
\noindent\begin{minipage}{\textwidth}\ProblemHeading{习题}{7.103}
求初值问题 $y^{\prime\prime}+y\sin x=0$、$y(0)=1$、$y^{\prime}(0)=0$ 的幂级数解，写至 $x^5$ 项。
\end{minipage}\par\vfill
\noindent\begin{minipage}{\textwidth}\ProblemHeading{习题}{7.104}
利用 Stirling 公式求下列极限：
\SubQ{1}{$\displaystyle\lim_{n\to\infty}\sqrt[n^2]{n!}$。}
\SubQ{2}{$\displaystyle\lim_{n\to\infty}\frac n{\sqrt[n]{n!}}$。}
\end{minipage}\par\vfill
\newpage
\noindent\begin{minipage}{\textwidth}\ProblemHeading{习题}{7.105}
研究下列级数的敛散性：
\SubQ{1}{$\displaystyle\sum_{n=3}^{\infty}\frac1{\ln(n!)}$。}
\SubQ{2}{$\displaystyle\sum_{n=1}^{\infty}\frac{n!e^n}{n^{n+p}}$，$p\in\mathbb R$。}
\end{minipage}\par\vfill
\noindent\begin{minipage}{\textwidth}\ProblemHeading{习题}{7.106}
证明：当 $n\to\infty$ 时，$\ln(n!)\sim\ln(n^n)$。
\end{minipage}\par\vfill
\noindent\begin{minipage}{\textwidth}\ProblemHeading{习题}{7.107}
求级数 $\displaystyle\sum_{n=1}^{\infty}\frac1{n(n+1)}\left(1+\frac12+\cdots+\frac1n\right)$ 的和。
\end{minipage}\par\vfill
\noindent\begin{minipage}{\textwidth}\ProblemHeading{习题}{7.108}
证明 $\displaystyle\sum_{n=0}^{\infty}(-1)^n\frac{2n+3}{(n+1)(n+2)}=1$。
\end{minipage}\par\vfill
\newpage
\noindent\begin{minipage}{\textwidth}\ProblemHeading{习题}{7.109}
设 $a_n$ 是正的递增数列。证明 $\sum_{n=1}^{\infty}(a_{n+1}/a_n-1)$ 收敛当且仅当 $a_n$ 有界。
\end{minipage}\par\vfill
\noindent\begin{minipage}{\textwidth}\ProblemHeading{习题}{7.110}
设 $\alpha>0$，$a_n$ 是正的递增数列。证明 $\displaystyle\sum_{n=1}^{\infty}\frac{a_{n+1}-a_n}{a_{n+1}a_n^\alpha}$ 收敛。
\end{minipage}\par\vfill
\noindent\begin{minipage}{\textwidth}\ProblemHeading{习题}{7.111（校订）}
设 $\Phi$ 在 $(0,+\infty)$ 上为正的严格增函数，并约定 $\Phi(0)=0$。非负数列 $a_n,b_n,c_n$ 满足\[a_{n+1}\le a_n-b_n\Phi(a_n)+c_na_n,\qquad\sum_{n=1}^{\infty}b_n=+\infty,\qquad\sum_{n=1}^{\infty}c_n<+\infty.\]证明 $a_n\to0$。
\end{minipage}\par\vfill
\noindent\begin{minipage}{\textwidth}\ProblemHeading{习题}{7.112}
设正数列 $a_n$ 满足 $\sum_{n=1}^{\infty}1/a_n$ 收敛。证明存在常数 $M>0$，使\[\sum_{n=1}^{\infty}\frac n{a_1+\cdots+a_n}\le M\sum_{n=1}^{\infty}\frac1{a_n}.\]
\end{minipage}\par\vfill
\newpage
\noindent\begin{minipage}{\textwidth}\ProblemHeading{习题}{7.113}
设 $a_n$ 为严格递增的实数列，且对每个正整数 $n$ 有 $a_n\le n^2\ln n$。证明 $\displaystyle\sum_{n=1}^{\infty}\frac1{a_{n+1}-a_n}$ 发散。
\end{minipage}\par\vfill
\noindent\begin{minipage}{\textwidth}\ProblemHeading{习题}{7.114}
若 $\sum_{n=1}^{\infty}|a_{n+1}-a_n|$ 收敛，称数列 $a_n$ 具有有界变差。
\SubQ{1}{证明有界变差数列必定收敛。}
\SubQ{2}{构造一个发散级数 $\sum_{n=1}^{\infty}a_n$，其通项数列具有有界变差。}
\end{minipage}\par\vfill
\noindent\begin{minipage}{\textwidth}\ProblemHeading{习题}{7.115}
在 $[0,1]$ 上递归定义 $f_0(x)=1$，$f_n(x)=\sqrt{x f_{n-1}(x)}$。证明 $f_n$ 一致收敛于一个连续函数。
\end{minipage}\par\vfill
\noindent\begin{minipage}{\textwidth}\ProblemHeading{习题}{7.116}
递归定义连续可微函数 $f_n:[0,1)\to\mathbb R$：$f_1=1$，$f_{n+1}^{\prime}=f_nf_{n+1}$（$0<x<1$），$f_{n+1}(0)=1$。证明对每个 $x\in(0,1)$，$\lim_{n\to\infty}f_n(x)$ 存在，并求极限函数。
\end{minipage}\par\vfill
\newpage
\noindent\begin{minipage}{\textwidth}\ProblemHeading{习题}{7.117}
设 $a>0$，$f_0$ 在 $[0,a]$ 上连续，递归定义 $f_n(x)=\int_0^x f_{n-1}(u)\,du$。证明 $f_n$ 在 $[0,a]$ 上一致趋于零。
\end{minipage}\par\vfill
\noindent\begin{minipage}{\textwidth}\ProblemHeading{习题}{7.118}
利用二项式级数计算 $\sqrt2$，保留四位小数。
\end{minipage}\par\vfill
